% Proposta de TCC — Monografia em Sistemas de Informação II
% Gerado a partir de docs/proposta-tcc2.md em 2026-10-06. Compilar com pdfLaTeX.
\documentclass[12pt,a4paper,oneside]{article}

% ----- codificação, idioma e fonte -----
\usepackage[utf8]{inputenc}
\usepackage[T1]{fontenc}
\usepackage[brazil]{babel}
\usepackage{tgheros}                       % TeX Gyre Heros (equivalente ao Arial)
\renewcommand{\familydefault}{\sfdefault}
\usepackage{microtype}

% ----- página ABNT: 3 cm (superior e esquerda), 2 cm (inferior e direita) -----
\usepackage[top=3cm,left=3cm,bottom=2cm,right=2cm]{geometry}
\usepackage{setspace}
\onehalfspacing
\usepackage{indentfirst}
\setlength{\parindent}{1.25cm}
\setlength{\parskip}{0.4em}

% ----- seções numeradas no estilo "1. Introdução" -----
\usepackage{titlesec}
\titleformat{\section}{\normalsize\bfseries}{\thesection.}{0.6em}{}
\titlespacing*{\section}{0pt}{1.6em}{0.8em}

% ----- listas, tabelas, matemática -----
\usepackage{enumitem}
\setlist{itemsep=0.15em,topsep=0.3em,parsep=0pt,leftmargin=*}
\usepackage{booktabs,tabularx,array,float}
\usepackage{amsmath,amssymb}

% ----- referências numeradas com seção numerada ("6. Referências Bibliográficas") -----
\usepackage[numbib,nottoc]{tocbibind}
\addto\captionsbrazil{\renewcommand{\refname}{Referências Bibliográficas}}

% ----- links -----
\usepackage{xurl}
\usepackage[hidelinks,breaklinks]{hyperref}
\urlstyle{same}

\begin{document}

% ===================== CAPA =====================
\begin{titlepage}
  \centering
  \vspace*{1.5cm}
  {\bfseries UNIVERSIDADE FEDERAL DE MINAS GERAIS\par}
  \vspace{1.2cm}
  {\bfseries Análise estrutural das comunidades no Twitter brasileiro durante o ciclo eleitoral de 2022\par}
  \vfill
  {\bfseries
    Aluno: Vinícius Pinho Galvão\\
    Orientadora: Eliane Cristina de Freitas Rocha\\
    Tipo de pesquisa: Mista (científica e tecnológica)\par}
  \vspace{1.5cm}
  {\bfseries Belo Horizonte\\ Outubro de 2026\par}
  \vspace*{2.5cm}
\end{titlepage}

% ===================== CORPO =====================
\section{Introdução}

O segundo turno da eleição presidencial de 2022 terminou com a vitória de Lula sobre Jair Bolsonaro por margem estreita. Pouco mais de dois meses depois, em 8 de janeiro de 2023, apoiadores de Bolsonaro invadiram e depredaram as sedes dos três poderes em Brasília, contestando o resultado das urnas \cite{bastos2023}. Para boa parte do país, esses fatos representavam um ataque à democracia. Para outra parcela, eram um ato legítimo de protesto. Não havia vocabulário comum entre as duas leituras.

Que o Brasil está dividido é quase um consenso intuitivo. O que é menos óbvio, e mais interessante de investigar, é a \textbf{forma} dessa divisão. Dois grupos opostos podem trocar informação o tempo todo ou se fechar em ambientes paralelos que quase não se tocam, e esses dois cenários têm consequências diferentes. No Twitter, a experiência de cada usuário é moldada pelo que ele escolhe amplificar e pelo que os algoritmos lhe entregam \cite{pariser2012}, \cite{sunstein2017}. Assim, duas pessoas expostas ao mesmo evento podem ter contato com versões quase opostas dele. Interian e Rodrigues \cite{interian2023} observaram um isolamento acentuado da comunidade de direita no Twitter durante a eleição de 2022 e levantaram a hipótese de que ele tenha influenciado o 8 de janeiro. Bastos e Recuero \cite{bastos2023} acrescentam que a narrativa de fraude foi também oferecida pela própria liderança, ao longo da campanha. A hipótese deste trabalho é que a divergência foi construída ao longo do ciclo eleitoral, e não de uma vez no 8 de janeiro. Por isso são analisados quatro eventos que formam um arco de tensionamento institucional:

\begin{itemize}
\item a mobilização de 7 de setembro de 2022;
\item o caso Roberto Jefferson, em 23 de outubro, quando o ex-deputado, aliado de Bolsonaro, reagiu a tiros e granadas a uma ordem de prisão do STF \cite{stf2022}–\cite{poder2022};
\item o dia do segundo turno, 30 de outubro;
\item a invasão de 8 de janeiro de 2023.
\end{itemize}

\textbf{O que foi feito no MSI I.} O objetivo geral do TCC1 era caracterizar a estrutura das comunidades políticas no Twitter brasileiro em torno do 8 de janeiro e torná-la compreensível por uma visualização interativa. O trabalho partiu do dataset desidratado Tweet\_Eleições\_2022 \cite{alabarse2024}. Como esse dataset não registra o autor dos tweets retuitados, construiu-se uma rede de co-retweet: dois usuários se ligam pela coincidência do que retuitaram, com peso dado pelo índice de Jaccard, e as comunidades são detectadas pelo algoritmo de Leiden. Os resultados relatados foram:

\begin{itemize}
\item uma pipeline reutilizável aplicada aos quatro eventos, com análise de sensibilidade dos parâmetros;
\item a caracterização estrutural: nos três eventos de confronto a rede se organiza em dois super-polos, com uma comunidade quase lacrada (hermeticidade de 0,97 a 0,999), achado consistente com o isolamento de uma das comunidades descrito em \cite{interian2023}; o dia do segundo turno apresentou a menor segregação pelas métricas de partição (modularidade e assortatividade).
\end{itemize}

O relatório do TCC1 apontou o próprio limite: as comunidades eram objetos estruturais, ainda sem conteúdo. Entre a entrega do relatório e esta proposta, deu-se o primeiro passo dessa camada: os tweets mais retuitados de cada comunidade foram hidratados seletivamente, um leitor web em versão beta foi construído e uma leitura manual preliminar desse material foi registrada. Essa leitura é insumo das perguntas abaixo e da validação prevista na E3.

\textbf{Problema e perguntas.} A pergunta central se mantém: \emph{o Twitter brasileiro estava polarizado no ciclo de 2022, e qual é a forma dessa polarização?} Uma revisão inicial dos resultados do TCC1, feita sobre os grafos já gerados, mostrou que, quando se somam as comunidades de cada lado, os dois super-polos são herméticos em grau comparável (0,98 a 0,9999 nos três eventos de confronto) e que o polo subdividido é tão ou mais fechado que o monolítico. A diferença entre eles está na organização interna: um polo aparece monolítico e o outro subdividido, e o lado subdividido muda de evento para evento; no 7 de setembro há ainda um pequeno terceiro grupo, coeso e distinto dos dois polos. No dia do segundo turno, o mesmo cálculo mostra dois polos fechados de tamanhos muito desiguais (96,6\% e 3,4\% dos nós), a reler junto com o efeito do filtro por palavra-chave do evento. Essa medição é preliminar e será incorporada à pipeline na E1. O TCC2 desdobra a pergunta central em quatro:

\begin{itemize}
\item \textbf{(P1)} as subdivisões internas dos polos são comunidades distintas ou artefatos da detecção?
\item \textbf{(P2)} o que cada polo diz sobre cada evento?
\item \textbf{(P3)} o que um polo vê que o outro não vê?
\item \textbf{(P4)} os ambientes são distintos a ponto de constituírem realidades paralelas?
\end{itemize}

Há também um problema de comunicação. Análises desse tipo costumam chegar a pesquisadores de outras áreas como resultados estáticos e difíceis de interrogar, e o leitor web do TCC1 se mostrou confuso para quem não vem da computação. A tese tecnológica do trabalho é que essa análise deve ser \textbf{navegável}, para que pesquisadores de comunicação e de ciência política consigam entender o método e explorar os achados por conta própria.

\textbf{O que muda em relação ao plano do MSI I.} O relatório do TCC1 previa rotular as comunidades como esquerda e direita por um score ideológico de autores, medir a radicalização de um painel de usuários ao longo do arco e animar o 8 de janeiro hora a hora. Esses itens saem do escopo. A leitura preliminar mostrou que ``esquerda'' e ``direita'' descrevem mal os polos encontrados (um deles reúne candidatos de centro e a imprensa de referência), e por isso os rótulos passam a emergir do conteúdo e da leitura dos participantes do grupo focal, não de uma classificação prévia. O painel de radicalização e o controle temporal ficam como trabalhos futuros.

\textbf{Objetivo geral.} Caracterizar a forma e o conteúdo da polarização no Twitter brasileiro ao longo do arco eleitoral de 2022–2023 e disponibilizar essa análise numa aplicação interativa compreensível e explorável por pesquisadores de comunicação e ciência política.

\textbf{Objetivos específicos (MSI II):}

\begin{enumerate}
\item Revisar a caracterização estrutural no nível dos polos e testar a robustez das subdivisões internas.
\item Hidratar seletivamente os tweets que discriminam as subdivisões robustas e, se o orçamento permitir, um quinto evento.
\item Identificar os tweets-assinatura de cada polo e de cada comunidade.
\item Produzir sínteses narrativas por evento e comunidade, com apoio de modelo de linguagem, ancoradas nos tweets que as sustentam e validadas.
\item Evoluir o leitor web para uma ferramenta de comparação entre polos e comunidades.
\item Avaliar a ferramenta num grupo focal com pesquisadores da área e publicá-la.
\end{enumerate}


\section{Referencial Teórico}

\textbf{Bolhas e câmaras de eco.} Pariser \cite{pariser2012} chamou de \emph{filter bubble} o universo informacional individual que os filtros algorítmicos montam para cada usuário, de forma invisível e sem que ele escolha. A noção de câmara de eco, de Sunstein \cite{sunstein2017}, enfatiza outro mecanismo: a homofilia, isto é, a tendência de pessoas semelhantes se agruparem e, deliberando entre iguais, se polarizarem. A rede de co-retweet registra o resultado agregado dos dois mecanismos, ou seja, o que cada usuário amplificou a partir do que lhe chegou. Por isso, a câmara de eco é tratada aqui como propriedade estrutural mensurável, e não como hipótese sobre o que cada usuário viu. Isso responde à crítica de que o termo costuma ser usado sem definição operacional e de que o isolamento completo é raro \cite{recuero2021}.

\textbf{Polarização no Twitter brasileiro.} O grupo de Raquel Recuero documentou, no debate sobre a hidroxicloroquina na pandemia, dois clusters fortemente demarcados e com dietas informativas distintas: num deles a desinformação circulava associada a veículos hiperpartidários; no outro predominavam veículos jornalísticos e institucionais \cite{recuero2021}. Nessa literatura, ``polarização assimétrica'' designa diferenças de dieta informativa entre os polos \cite{recuero2021}. O sentido usado aqui é outro: assimetria de organização interna, que esse trabalho não mede. O estudo mais próximo deste é o de Interian e Rodrigues \cite{interian2023}, que, numa rede de retweets dirigida com mais de 15 milhões de tweets de 2022, mostraram que a comunidade de direita respondia por 44,1\% da modularidade, contra 13,4\% da esquerda. Estudos recentes passam da estrutura ao conteúdo. Hoffmann da Silva e Recuero \cite{hoffmann2025} mostraram que os dois polos usaram ``democracia'' com sentidos opostos em agosto de 2022, o que antecipa a pergunta ``o que cada polo diz''. Bastos e Recuero \cite{bastos2023} documentaram a escalada da narrativa de fraude entre as lideranças.

\textbf{Da rede de retweet ao co-retweet.} Conover et al. \cite{conover2011} mostraram que a rede de retweets do Twitter político americano se divide em dois grupos segregados, com base na premissa de que o retweet funciona, em geral, como endosso. A rede dirigida exige conhecer o autor de cada tweet, informação ausente no dataset \cite{alabarse2024}. A rede de co-retweet liga usuários pelo conjunto de tweets que retuitaram em comum. Ela segue a lógica do acoplamento bibliográfico de Kessler \cite{kessler1963}, em que dois documentos são próximos quando citam as mesmas referências; formalmente, é a projeção da matriz de incidência usuário $\times$ tweet sobre o lado dos usuários, $C = M \cdot M^{\mathsf{T}}$. A analogia é proposta deste trabalho. A normalização por Jaccard é uma escolha analítica, declarada como tal. As comunidades resultantes são chamadas aqui de \textbf{comunidades discursivas}: grupos identificados por padrões de amplificação, sem rótulo prévio de orientação.

\textbf{Robustez da detecção de comunidades.} Uma partição produzida pelo Leiden \cite{traag2019} não deve ser tomada como verdade, por três motivos bem estabelecidos na literatura de detecção de comunidades:

\begin{itemize}
\item a modularidade admite muitas partições quase ótimas e distintas entre si;
\item a modularidade tem um limite de resolução, que o Constant Potts Model não tem;
\item execuções estocásticas variam, e o \emph{consensus clustering} permite medir a robustez de cada comunidade.
\end{itemize}

Se esses testes divergirem, o modelo de blocos estocásticos aninhado, com seleção por descrição mínima, oferece um critério independente para avaliar se uma subdivisão é sustentada pelos dados.

\textbf{Métodos digitais e ferramentas existentes.} A pesquisa em comunicação sobre plataformas tende a métodos mistos, que combinam mapeamento quantitativo e leitura qualitativa, e a literatura brasileira caracteriza clusters pela leitura manual de seus conteúdos e atores centrais \cite{recuero2021}, \cite{hoffmann2025}. O TCC2 mantém a leitura humana como validação, feita no grupo focal. As ferramentas mais usadas na área resolvem problemas vizinhos:

\begin{table}[H]
\centering
\singlespacing\small
\caption{Ferramentas existentes e posicionamento deste trabalho}\label{tab:ferramentas}
\begin{tabularx}{\textwidth}{>{\raggedright\arraybackslash}p{2.4cm} X X}
\toprule
\textbf{Ferramenta} & \textbf{O que faz} & \textbf{Diferença em relação a este trabalho} \\
\midrule
Gephi\footnotemark & Visualização e análise exploratória de grafos & Genérica; não liga comunidades ao conteúdo amplificado \\
NodeXL\footnotemark & Coleta e análise de redes de mídias sociais & Parte de coleta própria e de redes de interação; não trabalha com dados desidratados \\
4CAT\footnotemark & Captura e análise de dados de plataformas & Infraestrutura para quem analisa; não publica uma análise para terceiros explorarem \\
Hoaxy\footnotemark & Visualização pública da difusão de conteúdo & Foco na difusão de itens, não em comunidades nem na comparação entre polos \\
\bottomrule
\end{tabularx}
\end{table}
\addtocounter{footnote}{-4}
\stepcounter{footnote}\footnotetext{\url{https://gephi.org}}
\stepcounter{footnote}\footnotetext{\url{https://www.smrfoundation.org/nodexl/}}
\stepcounter{footnote}\footnotetext{\url{https://4cat.nl}}
\stepcounter{footnote}\footnotetext{\url{https://hoaxy.osome.iu.edu}}

A aplicação proposta não compete com essas ferramentas como software de análise. Ela é a \textbf{publicação interativa de uma análise}, que leva o leitor do grafo ao conteúdo.


\section{Metodologia}

São seis etapas, em ordem de dependência. Cada uma é implementada como estágio testado da pipeline existente.

\textbf{E1. Revisão estrutural.}

\begin{itemize}
\item Calcular a hermeticidade por super-polo, ao lado da hermeticidade por comunidade.
\item Abandonar a modularidade e a assortatividade como métricas de comparação entre eventos, porque misturam abertura entre polos com subdivisão interna.
\item Para cada par de comunidades entrelaçadas, testar a robustez: várias seeds do Leiden (com coassignação e ARI/NMI), varredura de resolução com o CPM como contraprova, variação do limiar de backbone e, se necessário, o SBM aninhado.
\item Registrar o critério de ``subdivisão robusta'' antes da execução.
\item \emph{Entregável:} tabela estrutural revisada e um veredito por par.
\end{itemize}

\textbf{E2. Hidratação seletiva.}

\begin{itemize}
\item Com o saldo da API do X, hidratar primeiro os tweets discriminantes dos pares robustos, isto é, os mais concentrados em um dos lados.
\item Com o saldo restante, hidratar um quinto evento do mesmo dataset, o que também demonstra que a pipeline pode ser reaplicada.
\end{itemize}

\textbf{E3. Camada de conteúdo.}

\begin{itemize}
\item \emph{Tweets-assinatura:} para cada tweet, calcular a fração dos retweets vinda de cada polo e de cada comunidade, sobre o log completo, acima de um piso de volume.
\item \emph{Sínteses narrativas por modelo de linguagem:} temas, protagonistas, antagonistas e tese central de cada comunidade em cada evento. Cada afirmação cita os tweets que a sustentam, e nenhuma síntese recebe rótulo ideológico prévio.
\item \emph{Validação:} checagem automática das citações, comparação com a leitura manual já realizada (o que foi capturado, omitido ou inventado) e revisão humana.
\end{itemize}

\textbf{E4. Ferramenta.} Evolução do leitor web, com dados estáticos:

\begin{itemize}
\item comparação lado a lado de polos e comunidades;
\item tweets-assinatura;
\item sínteses com links para os tweets;
\item explicação acessível do método;
\item cuidado ético na exibição de tweets.
\end{itemize}

\textbf{E5. Grupo focal.}

\begin{itemize}
\item Três a cinco pesquisadores de comunicação e ciência política realizam tarefas abertas na ferramenta, falando em voz alta enquanto a usam, seguidas de discussão.
\item \emph{Saídas:} dificuldades de interpretação, rotulagem manual dos polos e funcionalidades sugeridas.
\item Participação com termo de consentimento. A necessidade de apreciação por comitê de ética será verificada.
\end{itemize}

\textbf{E6. Iteração e publicação.} Incorporar as sugestões viáveis, publicar a ferramenta e redigir o artigo no formato IEEE.

\textbf{Uso de IA.} O modelo de linguagem é componente do método (E3) e é validado como descrito. Assistentes de IA também apoiam o desenvolvimento de código, a pesquisa bibliográfica e a revisão do texto, sob supervisão do autor, com todas as referências conferidas nas fontes.


\section{Resultados Esperados}

\begin{itemize}
\item \textbf{Analíticos:}
\begin{itemize}
  \item a forma da polarização nos quatro eventos, revisada no nível dos polos e das subdivisões;
  \item para cada evento e comunidade, a narrativa amplificada e os tweets-assinatura, que respondem às perguntas P1 a P4;
  \item resultados negativos, como subdivisões que se revelem artefatos, também serão reportados.
\end{itemize}
\item \textbf{Tecnológicos:}
\begin{itemize}
  \item a pipeline estendida, documentada e testada;
  \item a aplicação web publicada;
  \item uma medida da qualidade das sínteses (ancoragem e concordância com a leitura manual).
\end{itemize}
\item \textbf{De validação:}
\begin{itemize}
  \item o relato do grupo focal, com o que os pesquisadores conseguiram interpretar, os rótulos atribuídos aos polos e as melhorias sugeridas;
  \item o artigo final.
\end{itemize}
\end{itemize}


\section{Etapas e Cronograma}

\begin{table}[H]
\centering
\singlespacing\small
\caption{Cronograma previsto, em semanas a partir da entrega da proposta}\label{tab:cronograma}
\begin{tabularx}{\textwidth}{>{\centering\arraybackslash}p{1.8cm} X}
\toprule
\textbf{Semana} & \textbf{Atividade} \\
\midrule
1 & Entrega da proposta; registro do critério de robustez \\
2–3 & E1: revisão estrutural e testes de robustez \\
3–4 & E2: hidratação seletiva e quinto evento \\
4–6 & E3: tweets-assinatura, sínteses e validação \\
4–5 & Protocolo do grupo focal, consentimento e recrutamento \\
5–7 & E4: ferramenta para o grupo focal \\
7 & Apresentação parcial (pitch) \\
8–9 & E5: sessões do grupo focal e análise \\
9–10 & E6: iteração e publicação da ferramenta \\
10–12 & Redação do artigo, pôster e apresentação final \\
\bottomrule
\end{tabularx}
\end{table}


% ===================== REFERÊNCIAS =====================
\begin{thebibliography}{99}
\singlespacing
\bibitem{bastos2023} M. Bastos and R. Recuero, ``The insurrectionist playbook: Jair Bolsonaro and the National Congress of Brazil,'' \emph{Social Media + Society}, vol. 9, no. 4, Art. no. 20563051231211881, 2023, doi: 10.1177/20563051231211881.

\bibitem{pariser2012} E. Pariser, \emph{O Filtro Invisível: O que a Internet Está Escondendo de Você}. Rio de Janeiro: Zahar, 2012.

\bibitem{sunstein2017} C. R. Sunstein, \emph{\#Republic: Divided Democracy in the Age of Social Media}. Princeton, NJ, USA: Princeton Univ. Press, 2017.

\bibitem{interian2023} R. Interian and F. A. Rodrigues, ``Group polarization, influence, and domination in online interaction networks: A case study of the 2022 Brazilian elections,'' arXiv:2303.16859, 2023.

\bibitem{stf2022} Supremo Tribunal Federal, ``Ministro Alexandre de Moraes revoga domiciliar e determina nova prisão de Roberto Jefferson,'' \emph{Notícias STF}, 24 out. 2022. [Online]. Disponível em: \url{https://noticias.stf.jus.br/postsnoticias/ministro-alexandre-de-moraes-revoga-domiciliar-e-determina-nova-prisao-de-roberto-jefferson/}

\bibitem{ap2022} B. Prado and M. Savarese, ``Brazil politician, Bolsonaro ally refuses arrest, injures police after firing rifle,'' The Associated Press, via \emph{Global News}, Oct. 23, 2022. [Online]. Disponível em: \url{https://globalnews.ca/news/9220995/brazil-politician-bolsonaro-ally-refuses-arrest}

\bibitem{poder2022} L. Mendes and T. Angelo, ``Jefferson deu 50 tiros de fuzil e jogou 3 granadas, diz PF,'' \emph{Poder360}, 24 out. 2022. [Online]. Disponível em: \url{https://www.poder360.com.br/justica/jefferson-deu-50-tiros-de-fuzil-e-jogou-3-granadas-diz-pf/}

\bibitem{alabarse2024} L. Alabarse dos Santos, R. Mendes de Araujo, L. J. da Silva, A. G. D. Correa, O. B. Coelho, and I. C. A. de Oliveira, \emph{Tweet\_Eleições\_2022: Um dataset de tweets durante as eleições presidenciais brasileiras de 2022} [Dataset], Zenodo, 2024, doi: 10.5281/zenodo.11206577.

\bibitem{recuero2021} R. Recuero, F. B. Soares, and G. Zago, ``Polarização, hiperpartidarismo e câmaras de eco: como circula a desinformação sobre COVID-19 no Twitter,'' \emph{Contracampo}, vol. 40, no. 1, 2021, doi: 10.22409/contracampo.v40i1.45611.

\bibitem{hoffmann2025} M. L. Hoffmann da Silva and R. Recuero, ``Polarized democracies: unveiling distinct perspectives on the concept of Democracy in Brazilian social media political conversations,'' \emph{Estudos em Comunicação}, no. 41, pp. 3–22, 2025, doi: 10.25768/1646-4974n41v1a01.

\bibitem{conover2011} M. D. Conover, J. Ratkiewicz, M. Francisco, B. Gonçalves, F. Menczer, and A. Flammini, ``Political polarization on Twitter,'' in \emph{Proc. 5th Int. AAAI Conf. Weblogs and Social Media (ICWSM)}, 2011, pp. 89–96, doi: 10.1609/icwsm.v5i1.14126.

\bibitem{kessler1963} M. M. Kessler, ``Bibliographic coupling between scientific papers,'' \emph{American Documentation}, vol. 14, no. 1, pp. 10–25, 1963, doi: 10.1002/asi.5090140103.

\bibitem{traag2019} V. A. Traag, L. Waltman, and N. J. van Eck, ``From Louvain to Leiden: guaranteeing well-connected communities,'' \emph{Sci. Rep.}, vol. 9, Art. no. 5233, 2019, doi: 10.1038/s41598-019-41695-z.

\end{thebibliography}

\end{document}
